% !TEX program = xelatex
% ============================================================
% Подготовка к защите РГР (вариант 16), этап 2:
% функция Грина и аналитическое решение.
% Номера формул берутся из отчёта: сначала собрать ../report/main.tex,
% затем этот файл (latexmk defense2.tex).
% ============================================================
\documentclass[a4paper,12pt]{article}

\usepackage{fontspec}
\setmainfont{Times New Roman}
\usepackage{polyglossia}
\setmainlanguage{russian}
\setotherlanguage{english}
\setmonofont{Courier New}
\newfontfamily\cyrillicfonttt{Courier New}

\usepackage[left=22mm,right=18mm,top=18mm,bottom=20mm]{geometry}
\usepackage{setspace}
\setstretch{1.12}
\setlength{\parindent}{0pt}
\setlength{\parskip}{4pt plus 1pt}
\setlength{\emergencystretch}{2em}

\usepackage{amsmath,amssymb}
\renewcommand{\le}{\leqslant}
\renewcommand{\ge}{\geqslant}
\usepackage{array,longtable,booktabs}
\usepackage{enumitem}
\setlist{itemsep=1pt,topsep=2pt,leftmargin=1.6em}
\usepackage{titlesec}
\titleformat{\section}{\normalfont\large\bfseries}{\thesection}{0.6em}{}
\titleformat{\subsection}{\normalfont\normalsize\bfseries}{\thesubsection}{0.6em}{}
\titlespacing*{\section}{0pt}{14pt}{6pt}
\titlespacing*{\subsection}{0pt}{10pt}{4pt}

% Номера формул отчёта: \rf{eq:solution} -> (57)
\usepackage{xr-hyper}
\usepackage[unicode,hidelinks]{hyperref}
\externaldocument[R-]{../report/main}
\newcommand{\rf}[1]{(\ref*{R-#1})}
\newcommand{\rs}[1]{\ref*{R-#1}}

\newcommand{\pd}[2]{\dfrac{\partial #1}{\partial #2}}
\newcommand{\pdd}[2]{\dfrac{\partial^2 #1}{\partial #2^2}}
\DeclareMathOperator{\Div}{div}
\DeclareMathOperator{\Grad}{grad}

\newcounter{qa}
\newcommand{\q}[1]{\par\medskip\refstepcounter{qa}%
  \textbf{\theqa.~#1}\par\nopagebreak[4]}
\newcommand{\fhead}[2]{\par\medskip\textbf{Формула #1 — #2.}\par\nopagebreak[4]}

\begin{document}

\begin{center}
  {\Large\bfseries Подготовка к защите РГР, вариант 16\par}
  \vspace{4pt}
  {\large «Управление распределёнными системами».\\
   Этап 2: функция Грина и аналитическое решение (раздел 3 отчёта)\par}
\end{center}

\textbf{Как пользоваться.} Номера формул в скобках — номера из отчёта.
Материалы этапа~1 (постановка) — в отдельном файле. Порядок подготовки:
1)~выучить «цепочку из шести шагов» (раздел~\ref{sec:chain}) и итоговую
формулу; 2)~рассказ (раздел~\ref{sec:story}) — своими словами;
3)~пройти все формулы раздела~3 отчёта (раздел~\ref{sec:formulas});
4)~научиться воспроизводить на бумаге выводы из
раздела~\ref{sec:derivations}; 5)~прогнать вопросы, закрывая ответы.

\tableofcontents

% ============================================================
\section{Что держать в голове: цепочка из шести шагов}
\label{sec:chain}

\begin{enumerate}
  \item \textbf{Что такое $G$.} $\Delta G=-\delta(M,M_0)$, $G=0$ на боковой
  поверхности, $\partial G/\partial z=0$ на торцах. Условия те же по роду,
  что у $T$, но однородные. Смысл: температура от точечного источника.
  \item \textbf{Формула Грина} с вырезанным шаром вокруг $M_0$ даёт
  $T(M_0)=-\iint_{S_{\text{б}}}\mu\,\partial G/\partial n\,dS$.
  \item \textbf{Разложение по $z$ и $\varphi$.} По $z$: $Z''+\lambda Z=0$,
  $Z'(0)=Z'(h)=0$ $\Rightarrow$ $\cos(\pi nz/h)$. По $\varphi$:
  $\cos m(\varphi-\varphi_0)$. $G=\frac{1}{2\pi h}\sum\sum
  \varepsilon_m\varepsilon_n\cos\cos\cos\cdot g_{mn}(r,r_0)$.
  \item \textbf{Радиальное уравнение}
  $(rg')'-(m^2/r+k_n^2r)g=-\delta(r-r_0)$: непрерывность, скачок
  $r_0[g']=-1$, слева $I_m$ (ограничена на оси), справа комбинация,
  равная нулю при $r=a$. Вронскиан $I_mK_m'-I_m'K_m=-1/x$.
  \item \textbf{Производная $G$ на $r=a$:}
  $-a\,\partial g_{mn}/\partial r=I_m(k_nr_0)/I_m(k_na)$. Интеграл по углу
  оставляет только $m=0$.
  \item \textbf{Коэффициенты:} $\int_0^h\mu\,dz=T_0(h-c)$,
  $\int_0^h\mu\cos k_nz\,dz=-\frac{T_0h}{\pi n}\sin\frac{\pi nc}{h}$.
\end{enumerate}

\textbf{Итоговая формула} \rf{eq:solution}:
\[
T(r,z)=T_0\Bigl(1-\frac ch\Bigr)-\frac{2T_0}{\pi}\sum_{n=1}^\infty
\frac1n\sin\frac{\pi nc}{h}\,\frac{I_0(\pi nr/h)}{I_0(\pi na/h)}\,
\cos\frac{\pi nz}{h}.
\]
Как её «прочитать»: первое слагаемое — средняя по высоте температура
боковой поверхности; ряд — отклонения от среднего. Каждая гармоника
$\cos(\pi nz/h)$ на боковой поверхности имеет амплитуду из ряда Фурье
функции $\mu$, а к оси затухает с множителем $I_0(\pi nr/h)/I_0(\pi na/h)<1$.
Чем больше $n$ (чем мельче деталь), тем быстрее затухание.

% ============================================================
\section{Рассказ на три-четыре минуты}
\label{sec:story}

\textbf{Метод.} Решение строится методом функции Грина. Функция Грина —
это решение той же задачи с точечным источником единичной мощности внутри
и однородными граничными условиями: на боковой поверхности $G=0$, на
торцах $\partial G/\partial z=0$. Физически — температура, которую создаёт
точечный источник, если боковую поверхность охлаждать до нуля, а торцы
изолировать.

\textbf{Как $G$ выражает решение.} Беру вторую формулу Грина для $T$ и
$G$. Около точки $M_0$ функция $G$ ведёт себя как $1/(4\pi R)$, поэтому
вырезаю маленький шар. Интеграл по его сфере даёт $T(M_0)$. На торцах обе
нормальные производные равны нулю. На боковой поверхности $G=0$, поэтому
неизвестная $\partial T/\partial n$ выпадает, а $T=\mu$ известна.
Получается $T(M_0)=-\iint\mu\,\partial G/\partial n\,dS$. Вот зачем
условия на $G$ выбираются именно такими.

\textbf{Как найти $G$.} По $z$ и $\varphi$ условия однородные, поэтому
раскладываю $G$ в ряд по собственным функциям. По $z$ задача
Штурма–Лиувилля с условиями Неймана даёт косинусы $\cos(\pi nz/h)$,
$n=0,1,\dots$; по углу — $\cos m(\varphi-\varphi_0)$. Дельта-функцию
раскладываю в такие же ряды. Для коэффициента $g_{mn}(r)$ получается
обыкновенное уравнение с $\delta(r-r_0)$ справа — модифицированное
уравнение Бесселя. Его решения — $I_m$ (конечна на оси) и $K_m$ (бесконечна
на оси). Строю функцию Грина этого уравнения: слева от $r_0$ — $I_m$,
справа — комбинация, равная нулю при $r=a$, в $r_0$ — непрерывность и
скачок производной. Постоянную нахожу через вронскиан.

\textbf{Решение.} Дифференцирую $G$ по $r$ на боковой поверхности,
подставляю в интеграл. Интеграл по углу убивает все гармоники с $m\ne0$ —
так проявляется осевая симметрия. Остаётся ряд по $n$ с интегралами
$\int\mu\cos$, которые для ступеньки берутся элементарно. Ответ —
формула \rf{eq:solution}.

\textbf{Проверка.} Каждое слагаемое удовлетворяет уравнению, условиям на
торцах и на оси. При $r=a$ ряд становится рядом Фурье функции $\mu$ — он
сходится к $\mu$, в точке скачка к $T_0/2$, около скачка эффект Гиббса.
Внутри ряд сходится экспоненциально быстро. Средняя по высоте температура
на любом радиусе равна $T_0(1-c/h)$ — это баланс тепла. Численно проверено:
интеграл с $\partial G/\partial n$ совпадает с суммой ряда.

% ============================================================
\section{Обозначения раздела 3}
\label{sec:notation}

\begin{longtable}{@{}>{\raggedright\arraybackslash}p{2.8cm}>{\raggedright\arraybackslash}p{10.9cm}>{\raggedright\arraybackslash}p{2.2cm}@{}}
\toprule
Символ & Смысл & Единицы \\
\midrule
\endhead
$M$, $M_0$ & точка наблюдения $(r,\varphi,z)$ и точка источника $(r_0,\varphi_0,z_0)$ & — \\
$G(M,M_0)$ & функция Грина & 1/м \\
$\delta(M,M_0)$ & трёхмерная дельта-функция; $\iiint\delta\,dV=1$ & 1/м$^3$ \\
$\delta(z-z_0)$ и т.\,п. & одномерные дельта-функции & 1/м, 1/рад \\
$R_{MM_0}$ & расстояние между точками $M$ и $M_0$ & м \\
$v$ & регулярная часть функции Грина, $G-1/(4\pi R)$ & 1/м \\
$W$ & мощность точечного источника (только в объяснении смысла $G$) & Вт \\
$u$, $w$ & произвольные функции во второй формуле Грина \rf{eq:green2} & — \\
$B_\eta$, $\sigma_\eta$, $\eta$ & вырезанный шар, его сфера и радиус; $\eta\to0$ & м \\
$\Omega_\eta$ & цилиндр без шара & — \\
$S_{\text{б}}$ & боковая поверхность цилиндра & м$^2$ \\
$n_M$ & внешняя нормаль в точке $M$ & — \\
$Z$, $\Phi$ & функции в задачах Штурма–Лиувилля по $z$ и по $\varphi$ ($\Phi$ здесь не младшие члены из этапа 1) & — \\
$\lambda$, $\lambda_n$ & спектральный параметр и собственные значения по $z$ & 1/м$^2$ \\
$k_n=\pi n/h$ & «волновое число» $n$-й гармоники по $z$ (не теплопроводность $k$!) & 1/м \\
$\nu=m^2$ & собственные значения по $\varphi$ & — \\
$n$, $m$ & номера гармоник по $z$ и по $\varphi$ & — \\
$\varepsilon_n$, $\varepsilon_m$ & множитель Неймана: 1 при номере 0, 2 при номере $\ge1$ & — \\
$\|Z_n\|^2$ & квадрат нормы собственной функции & м \\
$g_{mn}(r,r_0)$ & радиальная функция Грина гармоники $(m,n)$ & — \\
$r_<$, $r_>$ & $\min(r,r_0)$ и $\max(r,r_0)$ & м \\
$\tau$ & полуширина отрезка при выводе скачка; $\tau\to0$ & м \\
$x=k_nr$ & безразмерный аргумент функций Бесселя & — \\
$I_m$, $K_m$ & модифицированные функции Бесселя I и II рода (Макдональда) & — \\
$y_2$ & решение, равное нулю при $r=a$ & — \\
$A$, $B$, $C$ & постоянные при построении $g_{mn}$ & — \\
$P(r_0,z_0;z)$ & ядро \rf{eq:poisson-kernel}: вклад пояска на высоте $z$ & 1/м \\
$\beta=a/h$ & отношение радиуса к высоте & — \\
$N$ & число слагаемых частичной суммы & — \\
$d$ & расстояние от точки до боковой поверхности в оценке сходимости & м \\
$\gamma$ & требуемая относительная точность & — \\
$Q(r)$ & полный тепловой поток через соосную цилиндрическую поверхность & Вт \\
$\psi$ & полярный угол с вершиной в точке скачка $(a,c)$ & рад \\
\bottomrule
\end{longtable}

\textbf{Ловушка обозначений:} $k$ — теплопроводность, $k_n$ — $\pi n/h$;
$K_m$ — функция Макдональда; $\rho$ — плотность (поэтому радиус шара
обозначен $\eta$); $\varepsilon_n$ — не малое число, а 1 или 2.

% ============================================================
\section{Разбор формул раздела 3}
\label{sec:formulas}

\subsection{Раздел 3.1 «Функция Грина и интегральное представление решения»}

\fhead{\rf{eq:G-def}}{определение функции Грина}
Уравнение $\Delta G=-\delta$: точечный единичный источник в $M_0$ (знак
минус — потому что в стационарном уравнении $k\Delta T=-F$, источник
положителен). Условия: $G=0$ на $r=a$ (I~род, однородное), $G_z=0$ на
торцах (II~род, однородное). Переменная — $M$, а $M_0$ — параметр.
Дополнительно: периодичность по $\varphi$, ограниченность вне $M_0$.
На оси условия нет: в трёхмерной задаче ось — внутренние точки.

\fhead{\rf{eq:G-sing}}{смысл равенства $\Delta G=-\delta$}
$G$ отличается от $1/(4\pi R)$ на гармоническую функцию $v$. Почему
$1/(4\pi R)$: $\Delta(1/R)=0$ при $R\ne0$ (проверка в сферических
координатах: $\frac{1}{R^2}(R^2\cdot(-1/R^2))'=0$), а поток
$\Grad(1/(4\pi R))=-\vec R/(4\pi R^3)$ через сферу радиуса $R$ равен
$-\frac{1}{4\pi R^2}\cdot4\pi R^2=-1$. По формуле Остроградского это и
значит «$\Delta=-\delta$»: интеграл от $\Delta$ по любой области вокруг
$M_0$ равен $-1$. Смысл $v$: поправка, обеспечивающая граничные условия
(аналог «изображения» источника).

\fhead{\rf{eq:green2}}{вторая формула Грина}
Из формулы Остроградского для вектора $u\Grad w-w\Grad u$:
$\Div(u\Grad w)=\Grad u\cdot\Grad w+u\Delta w$, то же с перестановкой,
разность: $u\Delta w-w\Delta u$ (смешанные слагаемые сокращаются).
Поток вектора через поверхность: $u\,\partial w/\partial n-w\,\partial u/\partial n$.

\fhead{\rf{eq:green-eps}}{формула Грина в области без шара}
В $\Omega_\eta$ обе функции гармоничны, левая часть — ноль. Граница
$\Omega_\eta$ — поверхность цилиндра плюс сфера $\sigma_\eta$.

\fhead{\rf{eq:sphere-limit}}{предел на маленькой сфере}
Нормаль внешняя к $\Omega_\eta$ — смотрит к $M_0$, то есть против $R$:
$\partial/\partial n=-\partial/\partial R$. $\partial(1/(4\pi R))/\partial n
=+1/(4\pi R^2)$. Интеграл: $T\approx T(M_0)$, умножаем на
$\frac{1}{4\pi\eta^2}\cdot4\pi\eta^2=1$. Второй интеграл:
$G\sim1/\eta$, площадь $\sim\eta^2$ — стремится к нулю.

\fhead{\rf{eq:T-green}}{решение через функцию Грина}
Знак минус — из переноса: $0=\iint_{S_\Omega}(\dots)+T(M_0)$. На торцах
$T_n=0$ и $G_n=0$ — ноль; на боковой $G=0$ и $T=\mu$. Нормаль на боковой
поверхности — по $r$; $dS=a\,d\varphi\,dz$ (дуга $a\,d\varphi$ на высоту
$dz$). Ключевая мысль: на каждой части границы известна одна величина из
двух; условие на $G$ «убивает» неизвестную.

\textbf{Тор вокруг окружности разрыва} (абзац после \rf{eq:T-green}).
На окружности $r=a$, $z=c$ $T$ разрывна, её производная там неограничена
($\sim1/\eta$). Вырезаем тор радиуса $\eta$: на его поверхности $T$
ограничена, $G\sim\eta$ (на границе $G=0$), $G_n$ ограничена, значит,
подынтегральное выражение ограничено, а площадь тора
$2\pi a\cdot2\pi\eta\to0$.

\textbf{Симметрия $G(M,M_0)=G(M_0,M)$.} Формула Грина для
$u=G(M,M_1)$, $w=G(M,M_2)$ с вырезанными шарами вокруг $M_1$, $M_2$:
граничные слагаемые на поверхности цилиндра равны нулю (оба удовлетворяют
одинаковым однородным условиям), сферы дают $G(M_1,M_2)$ и
$G(M_2,M_1)$ с разными знаками. Физически — принцип взаимности:
источник в $M_1$ нагревает $M_2$ так же, как источник в $M_2$ нагревает
$M_1$.

\subsection{Раздел 3.2 «Собственные функции по переменным $z$ и $\varphi$»}

\fhead{\rf{eq:SL-z}}{задача Штурма–Лиувилля по $z$}
Из условий Неймана на торцах. Почему по $z$: условия однородные именно
там. Перебор $\lambda<0$, $\lambda=0$, $\lambda>0$ — в отчёте; ответ:
$\lambda\ge0$ всегда (это видно и из тождества
$\lambda\int Z^2=\int Z'^2$, если умножить уравнение на $Z$ и
проинтегрировать по частям).

\fhead{\rf{eq:eigen-z}}{собственные значения и функции}
$k_n=\pi n/h$, $Z_n=\cos k_nz$, $n=0$ входит (постоянная) — в отличие от
условий Дирихле, где $\sin$ и $n\ge1$. Это важно: $n=0$ даёт среднюю
температуру $T_0(1-c/h)$.

\fhead{\rf{eq:norm-z}}{нормы}
$\int_0^h\cos^2(\pi nz/h)dz=h/2$ при $n\ge1$ ($\cos^2=\frac{1+\cos2x}{2}$),
$=h$ при $n=0$. Отсюда $\varepsilon_n$.

\fhead{\rf{eq:delta-z}}{разложение $\delta(z-z_0)$}
Коэффициент Фурье: $\frac{1}{\|Z_n\|^2}\int_0^h\delta(z-z_0)\cos k_nz\,dz
=\frac{\varepsilon_n}{h}\cos k_nz_0$ (фильтрующее свойство дельта-функции).

\fhead{\rf{eq:delta-phi}}{разложение $\delta(\varphi-\varphi_0)$}
Ряд Фурье на $[0,2\pi]$: коэффициенты $\frac{1}{\pi}\cos m\varphi_0$,
$\frac{1}{\pi}\sin m\varphi_0$, свободный член $\frac{1}{2\pi}$; синусы и
косинусы собираются в $\cos m(\varphi-\varphi_0)$.

\fhead{\rf{eq:delta-cyl}}{дельта-функция в цилиндрических координатах}
Множитель $1/r$ из-за $dV=r\,dr\,d\varphi\,dz$:
$\iiint\frac1r\delta\delta\delta\,r\,dr\,d\varphi\,dz=1$.

\subsection{Раздел 3.3 «Ряд для функции Грина»}

\fhead{\rf{eq:G-ansatz}}{вид ряда}
Множитель $\frac{\varepsilon_m\varepsilon_n}{2\pi h}$ взят из разложений
\rf{eq:delta-z}, \rf{eq:delta-phi}, чтобы уравнение для $g_{mn}$ имело
правую часть ровно $-\delta(r-r_0)$. Каждое слагаемое уже удовлетворяет
условиям на торцах и периодичности.

\fhead{\rf{eq:radial}}{уравнение для $g_{mn}$}
Подставили ряд в $\Delta G=-\delta$, сравнили коэффициенты, умножили на
$r$. Знак $+k_n^2r$: вторая производная косинуса по $z$ даёт $-k_n^2$,
переносим. $m^2/r$ — от $\frac{1}{r^2}\partial^2_\varphi$, умноженного
на $r$.

\fhead{\rf{eq:radial-bc}}{условия для $g_{mn}$}
$g(a)=0$ — из $G=0$ на $r=a$; ограниченность на оси — так как в 3D ось
внутри тела, $G$ там конечна (если $r_0>0$).

\fhead{\rf{eq:jump}}{условия сшивания}
Непрерывность — иначе $g''$ содержала бы $\delta'$. Скачок: интегрируем
уравнение по $[r_0-\tau,r_0+\tau]$:
$\bigl[rg'\bigr]_{r_0-\tau}^{r_0+\tau}-\int(\dots)g\,dr=-1$; интеграл
от ограниченной функции по отрезку длины $2\tau$ стремится к нулю.

\subsection{Раздел 3.4 «Радиальные функции Грина»}

\fhead{\rf{eq:mod-bessel}}{модифицированное уравнение Бесселя}
Подстановка $x=k_nr$: $\frac{d}{dr}=k_n\frac{d}{dx}$, $r\,g'=x\,y'$,
$(rg')'=k_n(xy')'$; уравнение $k_n(xy')'-k_n(m^2/x+x)y=0$ делим на $k_n$
и умножаем на $x$: $x(xy')'-(m^2+x^2)y=0$, то есть
$x^2y''+xy'-(x^2+m^2)y=0$. Обычное
уравнение Бесселя: $+(x^2-m^2)$ — его решения $J_m$ колеблются.
У нас знак другой: $I_m$, $K_m$ монотонны.

\fhead{\rf{eq:I-series}}{ряд для $I_m$}
Такой же, как для $J_m$, но без $(-1)^j$. Все члены положительны.
$I_0(x)=1+\frac{x^2}{4}+\frac{x^4}{64}+\dots$

\textbf{Свойства 2–3.} $K_m$ — второе решение, сингулярное в нуле, убывает
на бесконечности. Асимптотики: $I_m\sim e^x/\sqrt{2\pi x}$,
$K_m\sim\sqrt{\pi/(2x)}e^{-x}$ — «растущая» и «убывающая» экспоненты,
как $\operatorname{ch}$ и $e^{-x}$ для уравнения $y''-y=0$.

\fhead{\rf{eq:wronskian}}{вронскиан}
Формула Лиувилля: для $(py')'+qy=0$ произведение $pW$ постоянно. Здесь
$p=x$, поэтому $W=\text{const}/x$. Постоянная — по асимптотикам:
$W\approx I(-K)-IK=-2IK=-2\frac{e^x}{\sqrt{2\pi x}}\sqrt{\frac{\pi}{2x}}e^{-x}
=-\frac1x$.

\fhead{\rf{eq:g-n}}{$g_{mn}$ при $n\ge1$}
Слева — $I_m(k_nr)$ (ограничена). Справа — $y_2$, равная нулю при $r=a$.
Сшивание даёт $C=I_m(k_na)$. Проверка: при $r=a$ скобка $=0$; при
$r=r_0$ обе ветви совпадают; $g>0$.

\fhead{\rf{eq:g-0}}{$g_{m0}$ при $n=0$}
Уравнение Эйлера, решения $r^{\pm m}$ или $1,\ln r$. Для $m=0$:
$g=\ln(a/r_0)$ при $r<r_0$, $g=\ln(a/r)$ при $r>r_0$. Это функция
Грина «двумерного» кругового случая (как потенциал заряженной нити).

\subsection{Раздел 3.5 «Функция Грина задачи»}

\fhead{\rf{eq:G}}{итоговая функция Грина}
Двойной ряд. Свойства: симметрия, граничные условия, особенность
$1/(4\pi R)$ (численно: $4\pi RG\approx0{,}93$ при $R=0{,}04a$ —
недостающие $0{,}07$ дают $v$, отрицательная поправка от охлаждаемой
стенки), положительность (по принципу максимума для $G$ вне $M_0$:
на границе $G\ge0$, около $M_0$ $G\to+\infty$).

\subsection{Раздел 3.6 «Решение задачи»}

\fhead{\rf{eq:Gn}}{нормальная производная}
На $r=a$: $r_>=r$. Производная $y_2'(a)=k_n(I_m(k_na)K_m'(k_na)
-K_m(k_na)I_m'(k_na))=k_n\cdot(-1/(k_na))=-1/a$. Отсюда
$\partial g/\partial r=-I_m(k_nr_0)/(aI_m(k_na))$. Минус перед
$\partial G/\partial n$ делает всё положительным.

\textbf{Интеграл по углу} — $2\pi$ при $m=0$ и 0 при $m\ge1$: осевая
симметрия «автоматически».

\fhead{\rf{eq:T-mu}}{решение для любой $\mu(z)$}
$\frac{1}{h}\cdot\varepsilon_n\cdot$ (коэффициент затухания) $\cdot$
$\cos k_nz_0\cdot\int\mu\cos$. Если $\mu$ — любая функция от $z$,
формула та же.

\fhead{\rf{eq:poisson}, \rf{eq:poisson-kernel}}{ядро решения}
$T=\int\mu P\,dz$: температура — «взвешенная сумма» температур поясков
боковой поверхности. $P\ge0$, $\int_0^hP\,dz=1$ (слагаемое $n=0$), поэтому
$T$ — среднее значение $\mu$ с весом $P$ — отсюда снова
$\min\mu\le T\le\max\mu$.

\textbf{Коэффициенты} (перед \rf{eq:solution}):
$\int_c^h\cos\frac{\pi nz}{h}dz=\frac{h}{\pi n}\bigl[\sin\pi n
-\sin\frac{\pi nc}{h}\bigr]=-\frac{h}{\pi n}\sin\frac{\pi nc}{h}$.

\fhead{\rf{eq:solution}}{ответ}
Множитель $\frac{2T_0}{\pi}$: $\frac{1}{h}\cdot\varepsilon_n\cdot
\frac{T_0h}{\pi n}=\frac{2T_0}{\pi n}$. Знак минус — из
$-\sin(\pi nc/h)$. Условие $(r,z)\ne(a,c)$ — там решение не определено
(разрыв).

\fhead{\rf{eq:solution-dimless}}{безразмерная форма}
$\pi nr/h=\pi n\beta\xi$, $\pi na/h=\pi n\beta$. Видно: только $\beta=a/h$
и $c/h$.

\textbf{Абзац о разделении переменных.} Тот же ряд даёт метод Фурье:
$T=R(r)Z(z)$, $Z_n=\cos k_nz$, $R''+R'/r-k_n^2R=0$, ограниченное решение
$I_0(k_nr)$, коэффициенты из ряда Фурье $\mu$ при $r=a$. Функция Грина
даёт больше: формулу для любой $\mu$, ядро и симметрию.

\subsection{Раздел 3.7 «Проверка решения и его свойства»}

\fhead{\rf{eq:ratio-asympt}}{затухание гармоник}
Из асимптотик $I_0$:
$\frac{I_0(k_nr)}{I_0(k_na)}\approx\frac{e^{k_nr}/\sqrt{k_nr}}{e^{k_na}/\sqrt{k_na}}
=\sqrt{a/r}\,e^{-k_n(a-r)}$. При $r=0$: $1/I_0(x)\approx\sqrt{2\pi x}e^{-x}$.
Мажоранта — геометрическая прогрессия, отсюда равномерная сходимость
и почленное дифференцирование.

\textbf{Уравнение.} $I_0(k_nr)\cos k_nz$: $\partial_r^2+\frac1r\partial_r$
даёт $k_n^2(I_0''+I_0'/x)=k_n^2I_0$, $\partial_z^2$ даёт $-k_n^2$; сумма 0.

\fhead{\rf{eq:T-axis}}{температура на оси}
$I_0(0)=1$. Производная по $r$ на оси равна нулю: $I_0'=I_1$, $I_1(0)=0$.

\fhead{\rf{eq:mu-fourier}}{ряд Фурье граничной функции}
При $r=a$ отношение равно 1 — получаем разложение $\mu$ по косинусам.
Теорема Дирихле: для кусочно-гладкой функции ряд сходится к
$\frac{\mu(z-0)+\mu(z+0)}{2}$. В $z=c$: $\frac{0+T_0}{2}$.
Гиббс: выброс $\approx0{,}0895$ от величины скачка у любой частичной
суммы — ширина выброса уменьшается, высота нет (рисунок~3).

\fhead{\rf{eq:mean}}{средняя по высоте}
$\int_0^h\cos(\pi nz/h)\,dz=0$ при $n\ge1$. Смысл: поток через соосный
цилиндр $Q=-2\pi rk\int T_r\,dz=-2\pi rk\frac{d}{dr}\int T\,dz=0$ —
в стационаре и без источников внутри соосного цилиндра суммарный поток
через его поверхность ноль (торцы изолированы).

\textbf{Симметрия при $c=h/2$:} $T(r,h-z)=T_0-T(r,z)$ — замена
«$z\to h-z$, холодное $\leftrightarrow$ горячее».

\textbf{Предельные случаи.} $a\gg h$ — диск: далеко от края
$T\approx T_0(1-c/h)$. $a\ll h$ — стержень: $T\approx\mu(z)$, переход на
длине $\sim a$.

\textbf{Число слагаемых} $N\approx\frac{h}{\pi(a-r)}\ln\frac1\gamma$:
из условия $e^{-\pi N(a-r)/h}\approx\gamma$.

\textbf{Таблица 1.} Для $r/a\le0{,}5$ хватает 5–10 членов; для $r/a=0{,}9$
нужно около 50.

\subsection{Раздел 3.8 «Распределение температуры» и рисунки}

\textbf{Рисунок 3} — частичные суммы на $r=a$: $N=3$ — грубое
приближение, $N=60$ — почти ступенька с «ушами» Гиббса.

\textbf{Рисунок 4} — изотермы в осевом сечении. Цвет — $T/T_0$; левая
половина — зеркальная копия. Все изотермы сходятся в точку скачка
$(a,c)$; к торцам подходят под прямым углом (условие Неймана). Около
$(a,c)$ $T\approx T_0(\frac12+\psi/\pi)$ — как у плоской задачи с
разрывом: функция угла $\psi$ гармонична ($\Delta\psi=0$ на плоскости),
на верхней части границы ($\psi=\pi/2$) даёт $T_0$, на нижней
($\psi=-\pi/2$) — 0.

\textbf{Рисунок 5} — профили $T(z)$ на разных $r$. Чем ближе к оси, тем
плавнее. Все пересекают уровень $0{,}5$ около $z=c$, средние равны
$0{,}6$.

\textbf{Почему $T(0,c)>T_0/2$:} горячий пояс $h-c=0{,}6h$ длиннее
холодного $0{,}4h$, средняя температура $0{,}6\,T_0$.

% ============================================================
\section{Выводы, которые могут попросить воспроизвести}
\label{sec:derivations}

\subsection{Представление \texorpdfstring{$T(M_0)=-\iint\mu\,G_n\,dS$}{T(M0)}}
\label{der:rep}
1)~Вторая формула Грина. 2)~Вырезать шар $B_\eta$; левая часть 0.
3)~На сфере: $G_n=1/(4\pi\eta^2)$, интеграл $\to T(M_0)$; $GT_n$-часть
$\to0$. 4)~На торцах $T_n=G_n=0$, на боковой $G=0$, $T=\mu$.
5)~Перенести — знак минус.

\subsection{Задача Штурма–Лиувилля по $z$}
\label{der:sl}
Три случая по знаку $\lambda$ (см. отчёт, \rf{eq:SL-z}--\rf{eq:eigen-z}).
Быстрый аргумент $\lambda\ge0$: $\int_0^hZ(Z''+\lambda Z)dz=0$,
$\int ZZ''=[ZZ']_0^h-\int Z'^2=-\int Z'^2$ $\Rightarrow$
$\lambda\int Z^2=\int Z'^2\ge0$.

\subsection{Радиальная функция Грина при $m=0$, $n=0$ (самая простая)}
\label{der:g00}
$(rg')'=-\delta(r-r_0)$. При $r\ne r_0$: $rg'=\text{const}$, $g=A+B\ln r$.
Слева ограниченность $\Rightarrow$ $g=A$. Справа $g(a)=0$ $\Rightarrow$
$g=B\ln(a/r)$. Непрерывность: $A=B\ln(a/r_0)$. Скачок:
$r_0\bigl(g'(r_0+0)-g'(r_0-0)\bigr)=r_0(-B/r_0-0)=-1$ $\Rightarrow$ $B=1$.
Итог $g_{00}=\ln(a/r_>)$.

\subsection{Радиальная функция Грина при $n\ge1$}
\label{der:gn}
$g=AI_m(k_nr)$ слева, $g=By_2(r)$ справа,
$y_2=I_m(k_na)K_m(k_nr)-K_m(k_na)I_m(k_nr)$. Непрерывность: $A=y_2(r_0)/C$,
$B=I_m(k_nr_0)/C$. Скачок: $\frac{r_0}{C}\bigl(I_my_2'-k_nI_m'y_2\bigr)=-1$.
В скобке слагаемые с $K_m(k_na)$ сокращаются, остаётся
$I_m(k_na)k_n(I_mK_m'-I_m'K_m)=I_m(k_na)k_n\cdot(-\frac{1}{k_nr_0})$.
$C=I_m(k_na)$.

\subsection{Вронскиан $I_mK_m'-I_m'K_m=-1/x$}
\label{der:w}
Формула Лиувилля: $(xW)'=0$ для уравнения $(xy')'-(x+m^2/x)y=0$
(доказательство: $y_1(py_2')'-y_2(py_1')'=0$, это $(p(y_1y_2'-y_2y_1'))'=0$).
Константа по асимптотикам: $-1$.

\subsection{От $G$ к ответу}
\label{der:final}
$\partial g_{mn}/\partial r|_{r=a}=-\frac{1}{a}\frac{I_m(k_nr_0)}{I_m(k_na)}$
(через вронскиан в точке $k_na$). Подставить в $T=-\iint\mu G_n\,a\,d\varphi\,dz$;
$\int_0^{2\pi}\cos m(\varphi-\varphi_0)d\varphi=2\pi\delta_{m0}$; множители:
$a\cdot2\pi\cdot\frac{1}{2\pi h}\cdot\frac1a=\frac1h$. Интегралы от $\mu$ —
элементарные.

\subsection{Проверка граничных условий в ответе}
\label{der:check}
Торцы: $\partial_z\cos(\pi nz/h)\propto\sin(\pi nz/h)=0$ при $z=0,h$.
Боковая: ряд Фурье $\mu$. Ось: $I_0(0)=1$ — конечно.

% ============================================================
\section{Вопросы и ответы}
\label{sec:qa}

\subsection{О функции Грина}

\q{Что такое функция Грина вашей задачи?}
Функция $G(M,M_0)$, которая по $M$ удовлетворяет $\Delta G=-\delta(M,M_0)$
в цилиндре, $G=0$ на боковой поверхности и $\partial G/\partial z=0$ на
торцах. Это температура от точечного источника единичной мощности (при
$k=1$) в точке $M_0$ при охлаждённой до нуля боковой поверхности и
изолированных торцах.

\q{Почему условия для $G$ именно такие?}
Чтобы в формуле Грина исчезли неизвестные величины. На боковой поверхности
известна $T$, но не $\partial T/\partial n$ — поэтому $G=0$ там. На торцах
известна $\partial T/\partial n=0$, но не $T$ — поэтому
$\partial G/\partial n=0$.

\q{Почему в уравнении для $G$ минус перед $\delta$?}
Это соглашение, принятое в [2]: тогда $G>0$ и $G\approx1/(4\pi R)$
около источника, а уравнение совпадает со стационарным уравнением
$k\Delta T=-F$ для источника $F=\delta$, $k=1$. С $+\delta$ изменился бы
знак $G$ и знак в формуле \rf{eq:T-green}.

\q{Что такое дельта-функция?}
Обобщённая функция: равна нулю вне точки $M_0$, интеграл по любой
области, содержащей $M_0$, равен 1; $\iiint f\delta\,dV=f(M_0)$.
Модель точечного источника.

\q{Откуда $1/(4\pi R)$?}
Это фундаментальное решение уравнения Лапласа в пространстве:
сферически-симметричное гармоническое при $R\ne0$, поток градиента через
любую сферу $-1$. Физически — температура от точечного источника в
бесконечной среде.

\q{Почему функция Грина симметрична, и что это значит физически?}
Из формулы Грина для двух функций Грина с разными источниками. Физически —
взаимность: источник в $M_1$ создаёт в $M_2$ ту же температуру, что
источник в $M_2$ в $M_1$.

\q{Зачем вырезать шар? Почему нельзя сразу подставить?}
Формула Грина требует гладких функций, а $G$ в $M_0$ бесконечна. В области
без шара всё гладко, а предел интеграла по сфере даёт $T(M_0)$.

\q{Почему интеграл по сфере стремится именно к $T(M_0)$?}
$\partial G/\partial n\approx\frac{1}{4\pi\eta^2}$ (нормаль к центру),
площадь $4\pi\eta^2$, $T$ непрерывна — среднее $T$ по сфере $\to T(M_0)$.

\q{А что с разрывом $\mu$ в точке $z=c$?}
Вырезается тонкий тор вокруг окружности разрыва; на нём подынтегральная
функция ограничена, площадь $\to0$. Формула остаётся верной.
Дополнительно проверено численно.

\q{Чем функция Грина лучше простого разделения переменных?}
Ответ в виде $T=\int\mu P\,dz$ годится для любой $\mu$ — меняется только
интеграл. Видны физический смысл (источник, взаимность) и свойства
(положительность ядра, $\int P=1$, отсюда $0\le T\le T_0$).

\q{Получится ли тот же ответ разделением переменных?}
Да. $T=R(r)Z(z)$, по $z$ — те же косинусы, по $r$ — $I_0(k_nr)$
(ограниченное на оси), при $n=0$ — постоянная (а $\ln r$ отбрасываем).
Коэффициенты — из ряда Фурье $\mu$ при $r=a$.

\subsection{О собственных функциях и рядах}

\q{Почему раскладываем по $z$ и по $\varphi$, а не по $r$?}
По ним однородные условия (Неймана на торцах, периодичность), и
собственные функции простые — косинусы. По $r$ после этого остаётся ОДУ,
которое решается явно.

\q{Почему $n=0$ есть, а в задаче Дирихле его нет?}
При условиях Неймана постоянная удовлетворяет $Z'(0)=Z'(h)=0$, и
$\lambda=0$ — собственное значение. При Дирихле постоянная должна быть
нулём.

\q{Что такое $\varepsilon_n$?}
Множитель Неймана: 1 при $n=0$, 2 при $n\ge1$. Появляется из-за того,
что $\|\cos0\|^2=h$, а $\|\cos k_nz\|^2=h/2$.

\q{Почему собственные функции ортогональны?}
Задача Штурма–Лиувилля самосопряжённая; для разных $\lambda$:
$\int_0^h\cos k_nz\cos k_lz\,dz=0$ при $n\ne l$ — можно проверить прямо
по формуле произведения косинусов.

\q{Почему в радиальном уравнении $\delta(r-r_0)$ без $1/r$?}
В $\delta(M,M_0)$ есть $1/r$, но уравнение умножили на $r$. Из-за этого
и скачок записан как $r_0[g']=-1$.

\subsection{О функциях Бесселя}

\q{Почему появились функции Бесселя?}
Уравнение по $r$ — $g''+\frac1rg'-(k^2+m^2/r^2)g=0$; заменой $x=kr$ оно
становится модифицированным уравнением Бесселя. Слагаемое $\frac1rg'$ —
от цилиндрической геометрии.

\q{Почему модифицированные, а не обычные $J_0$?}
Знак при $k^2$: по $z$ гармоники — косинусы (колеблются), значит, по $r$ —
экспоненты (растут/убывают). Сумма вторых производных равна нулю: если
по одной переменной $-k^2$, по другой $+k^2$. $J_0$ появились бы при
разложении по $r$ (если бы однородные условия были на боковой
поверхности).

\q{Что такое $I_0$, как её посчитать?}
$I_0(x)=\sum_j\frac{(x/2)^{2j}}{(j!)^2}=1+\frac{x^2}{4}+\frac{x^4}{64}+\dots$
Положительна, возрастает, $I_0(0)=1$, растёт как $e^x/\sqrt{2\pi x}$.
$I_0(1)\approx1{,}266$, $I_0(\pi)\approx7{,}38$.

\q{Почему $K_m$ отбрасываем при $r<r_0$?}
$K_m\to\infty$ при $r\to0$, а на оси температура конечна.

\q{Что такое вронскиан и зачем он?}
$W=y_1y_2'-y_1'y_2$ — определитель из решений и их производных. Он
нужен для постоянной $C$ в функции Грина ОДУ и для производной $G$ на
границе. Для этого уравнения $W(I_m,K_m)=-1/x$.

\q{Чему равно $I_0'(x)$?}
$I_1(x)$; при $x=0$ ноль — отсюда $\partial T/\partial r=0$ на оси.

\subsection{О решении}

\q{Объясните каждую часть итоговой формулы.}
$T_0(1-c/h)$ — среднее $\mu$ по высоте; $\frac{2T_0}{\pi n}\sin\frac{\pi
nc}{h}$ — коэффициент Фурье $\mu$ (с минусом);
$\frac{I_0(\pi nr/h)}{I_0(\pi na/h)}$ — затухание $n$-й гармоники к оси;
$\cos\frac{\pi nz}{h}$ — собственная функция по $z$.

\q{Почему в ответе нет $\varphi$, $k$, $C$, $\rho$?}
$\varphi$ — интеграл по углу оставил только $m=0$ (осевая симметрия).
$k$, $C$, $\rho$ — уравнение Лапласа их не содержит.

\q{Почему решение в точке $(a,c)$ не определено?}
Там скачок граничной функции; пределы с разных сторон разные (от 0 до
$T_0$ в зависимости от направления). Ряд в этой точке даёт $T_0/2$, но
это не предел решения.

\q{Сходится ли ряд? Как быстро?}
Внутри — экспоненциально: члены $\sim\frac1n\sqrt{a/r}e^{-\pi n(a-r)/h}$.
На боковой поверхности — как ряд Фурье разрывной функции, $\sim1/n$,
неравномерно, с эффектом Гиббса.

\q{Сколько членов ряда надо взять?}
$N\approx\frac{h}{\pi(a-r)}\ln\frac1\gamma$; около оси — 5–10, около стенки
— десятки и сотни (таблица~1).

\q{Что такое эффект Гиббса?}
Частичные суммы ряда Фурье разрывной функции вблизи разрыва превышают
скачок примерно на 9\,\% при любом $N$; область выброса сужается, но его
высота не уменьшается. Внутри цилиндра его нет — высокие гармоники
подавлены.

\q{Почему средняя по высоте температура постоянна?}
Тепловой поток через любой соосный цилиндр равен нулю (нет источников,
торцы изолированы), а он пропорционален $\frac{d}{dr}\int_0^hT\,dz$.

\q{Выполняется ли $0\le T\le T_0$ в вашем решении?}
Да: ядро $P\ge0$ и $\int P\,dz=1$, значит, $T$ — взвешенное среднее
значений $\mu\in[0,T_0]$. Видно и по рисункам.

\q{Что будет при $c=h/2$?}
Чётные гармоники исчезнут, $T(r,h-z)=T_0-T(r,z)$; изотерма $T_0/2$ —
горизонтальная прямая $z=h/2$.

\q{Что будет при $c\to0$ или $c\to h$?}
$c\to0$: вся боковая поверхность при $T_0$, $\sin(\pi nc/h)\to0$, $T\to T_0$.
$c\to h$: $1-c/h\to0$, $\sin(\pi n)=0$, $T\to0$. Как и должно быть.

\q{Как проверяли численно?}
Программа \texttt{code/analytic.py}: вронскиан и скачок $g_{mn}$; симметрия
$G$; разностная проверка уравнения и условий на торцах; ряд при $r=a$
против $\mu$; интеграл с $\partial G/\partial n$ против суммы ряда
($0{,}7617\,T_0$ в обоих случаях). Функции Бесселя — из библиотеки
scipy, это допустимо: на этапе 3 запрет относится к готовым решателям
разностной задачи.

\q{Почему изотермы подходят к торцам под прямым углом?}
$\partial T/\partial z=0$: градиент горизонтален, а изотерма
перпендикулярна градиенту, то есть вертикальна у торца.

\q{Почему все изотермы выходят из точки $(a,c)$?}
Каждое значение между 0 и $T_0$ на границе достигается только в точке
разрыва. Около неё решение — функция угла $\psi$, изотермы — лучи.

\subsection{«А что, если…»}

\q{Если бы торцы держали при $T=0$?}
По $z$ были бы синусы, $n\ge1$, и решение
$\sum b_n\frac{I_0(\pi nr/h)}{I_0(\pi na/h)}\sin\frac{\pi nz}{h}$ со своими $b_n$.

\q{Если бы $T$ была задана на торцах, а боковая поверхность изолирована?}
Однородные условия были бы по $r$: $R'(a)=0$, собственные функции
$J_0(\mu_kr/a)$, $J_1(\mu_k)=0$; по $z$ — $\operatorname{ch}$, $\operatorname{sh}$.

\q{Если бы $\mu$ зависела от угла?}
Остались бы гармоники $m\ne0$ с $I_m(k_nr)$; формула \rf{eq:T-green}
та же.

\q{Если в цилиндре есть источники $F$?}
Добавится объёмный интеграл: $T(M_0)=\frac1k\iiint GF\,dV
-\iint\mu G_n\,dS$.

% ============================================================
\section{Ловушки}

\begin{enumerate}
  \item Перепутать $k$ (теплопроводность) и $k_n=\pi n/h$.
  \item Потерять гармонику $n=0$ — без неё ответ не имеет средней
  температуры $T_0(1-c/h)$.
  \item Написать $J_0$ вместо $I_0$ (неправильный знак при $k^2$).
  \item Забыть $\varepsilon_n$: при $n\ge1$ множитель 2, при $n=0$ — 1.
  \item Знак в $T(M_0)=-\iint\mu\,\partial G/\partial n\,dS$: минус, а
  $\partial G/\partial n<0$ на стенке, так что всё положительно.
  \item Сказать, что ряд на боковой поверхности сходится равномерно.
  \item Сказать, что в $(a,c)$ температура равна $T_0/2$: решение там
  разрывно, $T_0/2$ — только значение ряда.
  \item Сказать, что $G$ должна быть ограничена на оси «как условие»: в 3D
  ось — внутренние точки, условия там нет; ограниченность $g_{mn}$ на оси
  — следствие.
  \item Путать $\rho$ (плотность) и $\eta$ (радиус шара).
  \item Говорить «$K_m$ — функция Бесселя второго рода»: это $Y_m$
  (Неймана). $K_m$ — модифицированная функция второго рода (Макдональда).
\end{enumerate}

% ============================================================
\section{Самопроверка}

Могу ли я без подсказок:
\begin{itemize}[label=$\square$]
  \item определить функцию Грина и объяснить каждое условие;
  \item вывести $T(M_0)=-\iint\mu G_n\,dS$ с вырезанием шара;
  \item решить задачу Штурма–Лиувилля по $z$ (три случая $\lambda$);
  \item разложить $\delta(z-z_0)$ и $\delta(\varphi-\varphi_0)$;
  \item получить уравнение для $g_{mn}$ и условия сшивания;
  \item построить $g_{00}$ и $g_{mn}$ при $n\ge1$, вывести $C=I_m(k_na)$;
  \item объяснить, почему $I_m$, а не $J_m$, и почему $K_m$ отброшена слева;
  \item получить вронскиан по формуле Лиувилля и асимптотикам;
  \item вычислить $\partial G/\partial n$ на стенке и интеграл по углу;
  \item посчитать $\int\mu\cos k_nz\,dz$;
  \item написать итоговую формулу и объяснить каждый множитель;
  \item проверить все условия задачи для ответа;
  \item объяснить Гиббса, $T_0/2$ в точке скачка, среднюю температуру;
  \item объяснить рисунки 3–5 и таблицу 1.
\end{itemize}

\end{document}
